\chapter{A day's work}\label{ch:workflow}
\begin{goals}
\textbf{In this chapter} you take one day at an outcrop from start to finish: the compass's signed log
opened and checked, the readings split into sets, each set described, the fold axis found four ways, a
figure for the report, the project saved and sent, and a routine that does it again next time.
\end{goals}

The other chapters describe one thing each. This one strings them together, the way a working day
does. The outcrop is small and easy to picture: beds folded about one axis, two joint sets cutting them,
and a few hinge lineations measured where they could be seen. That is 37 readings in all: 31 planes and
6 lines.

Follow along on your own computer. Every number below is what wuffle prints for this log, and every
step points to the chapter that covers it in full.

\section{The day's log}\label{sec:wf-log}

In the field the log comes from the compass: \ui{publish →} writes \file{wuffle-log.csv} to
\file{Downloads}, its signature \file{wuffle-log.csv.sig} beside it, and the outcrop photos to
\file{Pictures} (\cref{sec:publish}). To follow this chapter without a phone, copy the listing below
into a plain-text file and save it as \file{wuffle-log.csv}. It is printed small so that each row stays
on one line; copy it exactly as it is.

\begin{lstlisting}[basicstyle=\ttfamily\tiny\color{ink}]
type,dip_dir_or_trend,dip_or_plunge,strike,timestamp,lat,lon,acc_m,north,declination_deg,photo
plane,164,75,74,2026-11-20T13:02:10.000Z,-30.071200,-51.120300,4,true,-15.21,
plane,44,63,314,2026-11-20T13:04:21.000Z,-30.071220,-51.120270,5,true,-15.21,
plane,62,44,332,2026-11-20T13:06:32.000Z,-30.071240,-51.120240,6,true,-15.21,
line,115,23,,2026-11-20T13:08:43.000Z,-30.071260,-51.120210,4,true,-15.21,
line,124,28,,2026-11-20T13:10:54.000Z,-30.071280,-51.120180,5,true,-15.21,wuffle-005-2026-11-20T131054.jpg
plane,86,29,356,2026-11-20T13:13:05.000Z,-30.071300,-51.120150,6,true,-15.21,
plane,201,68,111,2026-11-20T13:15:16.000Z,-30.071320,-51.120120,4,true,-15.21,
plane,65,36,335,2026-11-20T13:17:27.000Z,-30.071340,-51.120090,5,true,-15.21,
plane,192,54,102,2026-11-20T13:19:38.000Z,-30.071360,-51.120060,6,true,-15.21,
plane,167,71,77,2026-11-20T13:21:49.000Z,-30.071380,-51.120030,4,true,-15.21,
plane,288,83,198,2026-11-20T13:24:00.000Z,-30.071400,-51.120000,5,true,-15.21,
plane,79,31,349,2026-11-20T13:26:11.000Z,-30.071420,-51.119970,6,true,-15.21,
plane,50,54,320,2026-11-20T13:28:22.000Z,-30.071440,-51.119940,4,true,-15.21,
line,124,28,,2026-11-20T13:30:33.000Z,-30.071460,-51.119910,5,true,-15.21,
plane,290,86,200,2026-11-20T13:32:44.000Z,-30.071480,-51.119880,6,true,-15.21,
plane,172,72,82,2026-11-20T13:34:55.000Z,-30.071500,-51.119850,4,true,-15.21,
plane,179,45,89,2026-11-20T13:37:06.000Z,-30.071520,-51.119820,5,true,-15.21,
plane,295,85,205,2026-11-20T13:39:17.000Z,-30.071540,-51.119790,6,true,-15.21,wuffle-018-2026-11-20T133917.jpg
plane,159,73,69,2026-11-20T13:41:28.000Z,-30.071560,-51.119760,4,true,-15.21,
plane,292,84,202,2026-11-20T13:43:39.000Z,-30.071580,-51.119730,5,true,-15.21,
plane,169,77,79,2026-11-20T13:45:50.000Z,-30.071600,-51.119700,6,true,-15.21,
plane,296,83,206,2026-11-20T13:48:01.000Z,-30.071620,-51.119670,4,true,-15.21,
plane,201,72,111,2026-11-20T13:50:12.000Z,-30.071640,-51.119640,5,true,-15.21,
plane,293,82,203,2026-11-20T13:52:23.000Z,-30.071660,-51.119610,6,true,-15.21,
plane,36,72,306,2026-11-20T13:54:34.000Z,-30.071680,-51.119580,4,true,-15.21,
plane,195,63,105,2026-11-20T13:56:45.000Z,-30.071700,-51.119550,5,true,-15.21,
plane,55,50,325,2026-11-20T13:58:56.000Z,-30.071720,-51.119520,6,true,-15.21,
line,123,23,,2026-11-20T14:01:07.000Z,-30.071740,-51.119490,4,true,-15.21,
plane,190,49,100,2026-11-20T14:03:18.000Z,-30.071760,-51.119460,5,true,-15.21,
line,116,28,,2026-11-20T14:05:29.000Z,-30.071780,-51.119430,6,true,-15.21,
plane,160,70,70,2026-11-20T14:07:40.000Z,-30.071800,-51.119400,4,true,-15.21,
plane,171,34,81,2026-11-20T14:09:51.000Z,-30.071820,-51.119370,5,true,-15.21,
line,116,21,,2026-11-20T14:12:02.000Z,-30.071840,-51.119340,6,true,-15.21,
plane,296,84,206,2026-11-20T14:14:13.000Z,-30.071860,-51.119310,4,true,-15.21,
plane,160,77,70,2026-11-20T14:16:24.000Z,-30.071880,-51.119280,5,true,-15.21,
plane,175,39,85,2026-11-20T14:18:35.000Z,-30.071900,-51.119250,6,true,-15.21,
plane,291,83,201,2026-11-20T14:20:46.000Z,-30.071920,-51.119220,4,true,-15.21,
\end{lstlisting}

Each row is one reading as the compass recorded it: its kind, its attitude (dip direction and dip for a
plane, trend and plunge for a line), the strike, the time, where you stood and how well the phone knew
it, the north it refers to, and the photo taken there, if any. \cref{sec:format-log} gives every column.
Two things to notice before opening it. The \col{north} column says \code{true} on every row: the phone
had a position all day, so the readings were corrected for a declination of −15.21° as they were
taken (\cref{sec:true-north}). And two rows name a photo.

\section{Opening it}\label{sec:wf-open}

Open the desk (\cref{sec:opening}) and choose \menu{File>Open data or a project…}. Select the log, its
\file{.sig} and the two photos together, in one go, or drop all of them on the window at once. wuffle
gathers them into one field log and says what it found in the status line at the foot of the window:

\begin{lstlisting}
opened field log wuffle-log.csv — signature verified — key 3032 58b6 adb2 2e2f,
strongbox, signed 2026-11-20T18:40:02.311Z · 2 photos
\end{lstlisting}

The key's fingerprint is yours, not this one: it is your phone's key, and it is the same on every log
that phone signs. Compare it once with the fingerprint your phone shows, and after that the line tells
you at a glance that the file is your phone's record, byte for byte (\cref{sec:field-logs}).

The copy you typed from the listing has no \file{.sig} and no photos, so for it the line says instead:

\begin{lstlisting}
opened field log wuffle-log.csv — unsigned (2 named in the log weren't in the drop)
\end{lstlisting}

That is the honest answer for a file nobody signed. Everything else in this chapter works the same.

\shot{wf-opened}{The log opened: its planes and lines as two datasets on \ui{Net 1}, the planes' table
below, and the signature's verdict in the status line.}

The log becomes two datasets, \code{field 2026-11-20 · planes} (31 rows) and
\code{field 2026-11-20 · lines} (6 rows), named for the day the readings were taken and drawn on
\ui{Net 1}. The file itself is kept as it came, under \file{/project/field/field 2026-11-20/}, with its
signature and a \file{photos/} folder. In the console:

\begin{lstlisting}
wuffle /project $ verify
field 2026-11-20  (field/field 2026-11-20/wuffle-log.csv): signature verified — key 3032 58b6 adb2 2e2f, strongbox, signed 2026-11-20T18:40:02.311Z
wuffle /project $ ls -l "/project/field/field 2026-11-20"
d        0  photos
-     3040  wuffle-log.csv
-      374  wuffle-log.csv.sig
\end{lstlisting}

The datasets are yours to edit; the original never changes. In the tree's \ui{data} section both have
the arrow of a derived dataset, and hovering one says where it came from: \ui{made by field log from
field/field 2026-11-20/wuffle-log.csv} (\cref{sec:lineage}).

\subsection{A first look}

Look at the net before doing anything to it. The poles of the planes fall into three groups. Some trail
in an arc from the north round to the west-southwest: a spread like that, along a great circle, is what
the poles of folded beds do. Eight sit in a tight knot near the east of the primitive, and seven in
another near the north. The six lines, in their own colour, sit just inside the east-southeast rim.

The table confirms what the net suggests. Its first column, \col{dip dir/dip}, is wuffle's reading of
each row as an attitude (\cref{sec:table}); the planes read cleanly, and the heading says
\ui{field 2026-11-20 · planes · planes · 31 rows} with nothing skipped.

\section{Splitting the sets}\label{sec:wf-split}

The three groups are three different things: bedding, and two joint sets. Mixed together their
statistics mean nothing: the \code{stats} of all 31 planes is a mean of \code{146/63} with
$\kappa$~2.8, a description of no surface in the outcrop. Each set needs to be a dataset of its own.

Split them with selections (\cref{ch:select}). Selections are drawn on the sphere, so a cone means the
same rows whatever the net's projection or rotation, and a selection made into a dataset keeps a record
of where it came from.

\begin{tryit}
\begin{enumerate}
\item Open the console (\key{`}) and type \code{select cone 292/84 15}. A cone about a plane is a cone
about its pole, so this takes every pole within 15° of the pole to \code{292/84}. The console answers
\ui{cone: 8 selected}, and the eight poles of the eastern knot turn violet inside a dashed outline.
\item Right-click the planes layer in the tree and choose \ui{Selection → new dataset}. A dataset called
\code{field 2026-11-20 · planes · selection} appears in the \ui{data} section. Right-click it, choose
\ui{Rename…} and call it \code{J1}.
\item Type \code{select cone 166/74 15} (\ui{cone: 7 selected}), make it a dataset the same way, and
call it \code{J2}.
\item For the bedding, take everything and then remove the joints: \code{select all}, then
\code{select - cone 292/84 15} and \code{select - cone 166/74 15}. The 16 planes left are the beds.
Make them a dataset called \code{bedding}.
\item Type \code{select none}.
\end{enumerate}
\end{tryit}

\shot{wf-split}{The first joint set selected by a cone typed in the console: eight poles violet inside
the dashed cone, and the console's answer.}

You can draw the same cones with the mouse instead, with the \ui{cone} tool above the net
(\cref{sec:shapes}); typing them gives round numbers you can write in your notes and repeat. In step~4,
\code{select all} also selects the six lines, which does no harm: \ui{Selection → new dataset} on the
planes layer takes only the planes' selected rows.

The \ui{data} section now lists \code{J1} (8 rows), \code{J2} (7) and \code{bedding} (16), each with
the lineage \ui{made by selection from field 2026-11-20 · planes} (\cref{sec:select-dataset}). Nothing
was copied out of the project or retyped: if you find a misread plane tomorrow and correct it in the
log's dataset, you know exactly which sets were cut from it.

\section{A net of the sets}\label{sec:wf-sets}

Give the sets a net of their own, and keep \ui{Net 1} as the raw day.

\begin{enumerate}
\item Choose \menu{Diagram>New net}. The new net starts with \ui{Net 1}'s two layers, as views of the
same data (\cref{sec:views}). Right-click it in the tree, choose \ui{Rename…} and call it \code{Sets}.
\item Remove the two layers it came with: right-click each and choose \ui{Remove from this diagram}.
The data stays.
\item From the \ui{data} section, draw \code{bedding}, \code{J1}, \code{J2} and the log's lines on
\ui{Sets}, in that order: right-click each and choose \menu{Draw on>Sets}, or drag it onto the net's
name in the tree.
\item Give each a colour you can tell apart in the layer's properties (double-click it,
\cref{sec:properties}). The figures here use blue for the bedding, orange for J1, green for J2 and brown
for the lines.
\end{enumerate}

Now choose what each layer draws (\cref{ch:elements}). The tree's checkboxes do it, or, in the console,
with \ui{Sets} active:

\begin{lstlisting}
use Sets
contour kamb 3 bedding
show v3 bedding
show mean J1
show mean J2
show mean field
\end{lstlisting}

The bedding is contoured (with 16 poles the contours are patches along the girdle rather than a band),
and its third eigenvector, V3, is drawn as a point. Each joint set and the lines get their Fisher mean with its cone of confidence. (\code{field} is
enough to name the lines layer: a layer is found by the start of its name, \cref{sec:commands}.)

\shot{wf-sets}{The \ui{Sets} net: the bedding contoured with its V3, the means and cones of J1, J2 and
the hinge lines; the statistics tab below shows a card per layer.}

\subsection{Reading the cards}

The \ui{statistics} tab has one card per visible layer (\cref{sec:stats-card}). For this outcrop:

\begin{center}\small
\begin{tabularx}{\linewidth}{@{}>{\raggedright}p{0.17\linewidth}>{\raggedright}p{0.22\linewidth}>{\raggedright}p{0.1\linewidth}>{\raggedright}p{0.1\linewidth}L@{}}\toprule
\sffamily\bfseries layer & \sffamily\bfseries mean & \sffamily\bfseries $\kappa$ & \sffamily\bfseries $\alpha_{95}$ & \sffamily\bfseries shape \\\midrule
bedding & \code{152/29} & 2.6 & 27.3° & K 0.01, C 6.80: a girdle; V3 \code{120→25} \\
J1 & \code{293/84} & 557.7 & 2.2° & K 3.72, C 7.73: a tight cluster \\
J2 & \code{164/74} & 174.1 & 4.2° & K 4.74, C 6.16: a cluster \\
lines & \code{120→25} & 201.6 & 4.2° & K 4.81, C 6.40: a cluster \\
\bottomrule\end{tabularx}
\end{center}

The joint sets are what the field suggested: steep, well defined, J1 dipping 84° toward 293 and J2
74° toward 164. Their cones of confidence are a few degrees across.

The bedding is the interesting card. Its mean, \code{152/29} with $\kappa$~2.6 and a cone of
confidence 27° wide, is not the attitude of any bed: it is the average of two limbs that dip in
opposite directions. The eigen analysis is what describes it. The eigenvalues are
\code{0.516 0.483 0.001}: two large and one almost nothing, which is a girdle, and Woodcock's
K~0.01 says the same \parencite{woodcock1977}. The poles lie on a great circle, \code{girdle 300/65},
whose pole V3 is \code{120→25}. That is the fold axis, and the girdle is the fold's profile plane
(\cref{sec:stats-eigen}).

\shot[0.7\linewidth]{wf-card}{The bedding's card: a girdle (S3 0.001, K 0.01), with V3 at \code{120→25}
and the girdle \code{300/65}.}

\section{Roses and a fabric plot}\label{sec:wf-rose}

With \ui{Sets} active, choose \menu{Diagram>New rose diagram}. Like a new net, the rose starts with the
active diagram's layers, and by default it bins strikes, which skips the lines by itself
(\cref{sec:rose}).

\shot[0.55\linewidth]{wf-rose}{A rose of the strikes of the bedding, J1 and J2 together, in 10° sectors.}

The longest pair of petals, toward 020 and 200, is J1, whose eight strikes all lie between 198 and 206.
J2's strikes, 069 to 082, make the petals toward the east-northeast and west-southwest. The bedding's
strikes are spread round the rest of the rose: the two limbs strike apart, and between them they fill
it. A rose is a good summary of the joints for a report, and a poor one of folded beds, which the net
shows better.

A fabric plot compares the shape of each set's distribution in one picture (\cref{sec:fabric}). With
\ui{Sets} active, choose \menu{Diagram>New fabric plot}, then right-click it in the tree and choose
\menu{Plot>Woodcock K–C}.

\shot[0.55\linewidth]{wf-fabric}{The fabric plot, Woodcock's form: the bedding far down to the right (a
girdle), J1, J2 and the hinge lines up to the left (clusters).}

The bedding sits at the bottom right: a girdle, with K close to zero. J1, J2 and the lines sit high on
the left, well above the line K~=~1: clusters, J1 the strongest. On the default Vollmer triangle all three
clusters sit at the P corner together, their labels stacked one under another; the Woodcock form spreads
them apart by strength, which is what a day like this wants to show.

\section{A figure of the sets}\label{sec:wf-figure}

The net of the sets is the figure the report wants. Make \ui{Sets} the active tab and choose
\menu{Figure>New figure from the active net} (\cref{sec:new-figure}).

\shot{wf-figure}{The figure from the \ui{Sets} net, open in its own tab: the title, the net, a legend, the
statistics table and a caption, each a frame you can move.}

The page has five frames: the title \code{Sets}; the net, drawn as vectors; a legend of the four
layers with their counts, \code{bedding (16)}, \code{J1 (8)}, \code{J2 (7)} and
\code{field 2026-11-20 · lines (6)}; a table of the statistics; and a caption that says what is drawn and
where each dataset came from:

\begin{quote}\small
Equal-area (Schmidt), lower hemisphere · Kamb contours at 2σ intervals (k = 3). Planes as dip
direction/dip, lines as trend→plunge. Derived: bedding: selection of field 2026-11-20 · planes; J1:
selection of field 2026-11-20 · planes; J2: selection of field 2026-11-20 · planes;
field 2026-11-20 · lines: field log of field/field 2026-11-20/wuffle-log.csv. Made with wuffle,
\textit{(today's date)}.
\end{quote}

The caption is written from the project, not typed: the lineage you built by selecting the sets is what
it reports. Click it to add a sentence of your own, such as where the outcrop is. Then
\ui{Export PDF} writes \file{Sets\_figure.pdf}, the page as a vector PDF (\cref{sec:figure-export}). If
you change the net later, \ui{↻ from the net} redraws the figure and keeps your layout.

\section{The fold axis, four ways}\label{sec:wf-axis}

The bedding's V3 says the fold axis is \code{120→25}. One estimate is a claim; several independent ones
that agree are a result. This log holds three more.

\paragraph{β.} On the bedding's card, press \ui{β → dataset}. wuffle intersects every pair of the 16
beds, 120 lines in all, and keeps them as a dataset, \code{bedding · β}, drawn on \ui{Sets}
(\cref{sec:derived}). Its own card gives their mean: \code{120→24}, $\kappa$~131.4.

\begin{note}
The 120 lines of a $\beta$ diagram come from only 16 measurements, so they are far from independent:
their $\alpha_{95}$ of 1.1° says how well the mean of those lines is known, not how well the axis is.
Read the spread of the points instead.
\end{note}

\paragraph{The limbs' intersection.} Split the bedding into its two limbs. On \ui{Sets}, choose the
\ui{lasso} tool and draw round the eight bedding poles in the west-southwest of the net; make them a
dataset from the bedding layer's menu and call it \code{NE limb}. Lasso the eight poles near the north
and call them \code{SW limb}. Draw both on \ui{Sets}. Their cards give the two limbs' mean planes,
\code{056/46} and \code{190/53}. On the \code{NE limb} card, choose \ui{SW limb} from \ui{∩ intersect
with…}: wuffle keeps the line where the two mean planes meet as \code{NE limb ∩ SW limb}, at
\code{120→24} (\cref{sec:stats-intersect}). Its table gives the angle between the two planes as well:
88.71°.

\paragraph{The hinge lines.} The six lineations measured on the hinge have a mean of \code{120→25},
$\kappa$~201.6, with a cone of confidence of 4.2°: the card of the lines layer already says so. Press
its \ui{mean → dataset} to keep it as a dataset with the others.

The same in the console, with \ui{Sets} active (and the limbs drawn there):

\begin{lstlisting}
wuffle /project $ beta bedding
→ bedding · β (120 rows), drawn on Sets
wuffle /project $ intersect NE SW
→ NE limb ∩ SW limb (1 row), drawn on Sets
wuffle /project $ mean field
→ field 2026-11-20 · lines · mean (1 row), drawn on Sets
\end{lstlisting}

\shot[0.8\linewidth]{wf-axis}{The fold axis four ways. The bedding poles lie on their best-fit girdle;
in the east-southeast, the grey $\beta$ points, the hinge lines with their mean and cone, the bedding's
V3 and the limbs' intersection all crowd round \code{120→25}. (The joint sets are hidden, and the limbs'
layers too, since their poles are the bedding's.)}

To draw that picture, untick \code{J1} and \code{J2} and the two limb layers, untick the bedding's
contours and tick its \ui{V3 great circle}, and give the $\beta$ layer a small grey symbol in its
properties so it doesn't hide the rest.

All four agree within a degree:

\begin{center}\small
\begin{tabularx}{\linewidth}{@{}>{\raggedright}p{0.36\linewidth}>{\raggedright}p{0.16\linewidth}L@{}}\toprule
\sffamily\bfseries estimate & \sffamily\bfseries axis & \sffamily\bfseries from \\\midrule
V3 of the bedding girdle & \code{120→25} & the 16 bedding poles, by their eigen analysis \\
mean of the $\beta$ lines & \code{120→24} & every pair of beds \\
the limbs' intersection & \code{120→24} & the two limbs' mean planes \\
the hinge lineations' mean & \code{120→25} & six lines measured on the hinge \\
\bottomrule\end{tabularx}
\end{center}

That agreement is the finding. The fold is cylindrical, its axis plunges 25° toward 120, and its limbs
stand at about 90° to each other. Three of the estimates come from the same beds by different
arithmetic, but the fourth comes from different measurements altogether, which is what makes it
worth having.

\section{Saving, and sending}\label{sec:wf-save}

First give the project a name: \menu{File>Rename project…}, and type \code{serra-2026-11-20}
(\cref{sec:project-name}). The tab's title becomes \ui{serra-2026-11-20 — wuffle}, and the files are
saved under that name.

\menu{File>Save project (.zip.wuffle)} (\key{Ctrl}+\key{S}) writes the whole day as one file,
\file{serra-2026-11-20.zip.wuffle} (\cref{sec:project-file}). Inside it are the log, byte for byte, with
its signature and photos under \file{field/}; every dataset as a plain CSV, the sets and the derived ones
with their lineage; both nets, the rose, the fabric plot and the figure. For this day it comes to about
55~KB.

To send the day to someone who doesn't have wuffle, choose \menu{File>Save as page (one .html, the
project inside)} (\cref{sec:save-page}). It writes \file{serra-2026-11-20.wuffle.html}, about 800~KB: wuffle
itself with the project inside. Whoever opens it lands on the desk with your nets, sets and figure in
front of them, offline, with nothing to install, and their changes are kept in their own browser, apart
from yours. The signed log travels inside it too, so they can run \code{verify} themselves.

\section{Doing it again}\label{sec:wf-routine}

Tomorrow brings another log from the same outcrop. The sets will need choosing again, since that takes
a look at the net, but the description of them is the same every time. Keep it as a routine
(\cref{sec:routines}): a script in \file{/home}, which is saved with the project.

\begin{tryit}
\begin{enumerate}
\item In the console, write the routine:
\begin{lstlisting}
printf 'use Sets\nfor s in bedding J1 J2 field; do stats $s; done\n' > /home/sets.sh
\end{lstlisting}
\item Run it, keeping what it prints in a report:
\begin{lstlisting}
. /home/sets.sh > /home/report.txt
\end{lstlisting}
\item Read the report back with \code{cat /home/report.txt}. It has 13 lines: the active diagram, then
the count, mean and eigenvalues of each set.
\end{enumerate}
\end{tryit}

\shot{wf-routine}{The routine written into \file{/home}, run into a report, and the first six lines of
the report read back.}

The report starts like this:

\begin{lstlisting}
● Sets
bedding · n 16
mean 152/29 · pole 332→61 · κ 2.6 · α95 27.3°
S 0.516 0.483 0.001 · K 0.01 · C 6.80
J1 · n 8
mean 293/84 · pole 113→06 · κ 557.7 · α95 2.2°
\end{lstlisting}

The routine and its report are files in \file{/home}, so they go into the next save and into a page you
send. Grow the routine as the work settles: \code{beta bedding} and \code{intersect NE SW} when the
limbs are drawn, \code{contour kamb 3 bedding}, \code{verify} at the top to put the log's signature in the
report. \cref{ch:console} has the commands and \cref{app:commands} every one of them.

\begin{why}
The selections that make the sets are judgement: where one set ends and the next begins is a call you
make looking at the net. What follows from them is arithmetic. Keeping the judgement in the open (each
set's lineage names its cone) and the arithmetic in a routine means tomorrow's numbers come out the same
way as today's, and anyone you send the project to can see both.
\end{why}
